%%%%%%%%%%%%%%%%%
% CV tailored for Mi4 - Data Integration Specialist
%%%%%%%%%%%%%%%%%

\documentclass[8pt,a4paper,ragged2e,withhyper]{altacv}

\geometry{left=1.25cm,right=1.25cm,top=.5cm,bottom=1.cm,columnsep=1.2cm}

\usepackage{paracol}
\usepackage{fontawesome5}
\usepackage{hyperref}

\iftutex 
  \setmainfont{Roboto Slab}
  \setsansfont{Lato}
  \renewcommand{\familydefault}{\sfdefault}
\else
  \usepackage[rm]{roboto}
  \usepackage[defaultsans]{lato}
  \renewcommand{\familydefault}{\sfdefault}
\fi

\definecolor{SlateGrey}{HTML}{2E2E2E}
\definecolor{LightGrey}{HTML}{666666}
\definecolor{DarkPastelRed}{HTML}{8AC0D9}
\definecolor{PastelRed}{HTML}{00B0DA}
\definecolor{GoldenEarth}{HTML}{B4AD0C}

\colorlet{name}{black}
\colorlet{tagline}{PastelRed}
\colorlet{heading}{DarkPastelRed}
\colorlet{headingrule}{GoldenEarth}
\colorlet{subheading}{PastelRed}
\colorlet{accent}{PastelRed}
\colorlet{emphasis}{SlateGrey}
\colorlet{body}{LightGrey}

\definecolor{violet}{HTML}{5A2582}
\definecolor{rouge}{HTML}{F53B3B}
\definecolor{noir}{HTML}{00232C}
\definecolor{bleu}{HTML}{00B0DA}
\definecolor{rose}{HTML}{F831A0}
\definecolor{jaune}{HTML}{FFEC00}
\definecolor{vert}{HTML}{34AD6C}

\renewcommand{\namefont}{\Huge\rmfamily\bfseries}
\renewcommand{\personalinfofont}{\footnotesize}
\renewcommand{\cvsectionfont}{\LARGE\rmfamily\bfseries}
\renewcommand{\cvsubsectionfont}{\large\bfseries}
\renewcommand{\cvItemMarker}{{\small\textbullet}}
\renewcommand{\cvRatingMarker}{\faCircle}

\input{pubs-num.tex}
\addbibresource{sample.bib}

\begin{document}

\name{BOILEAU Guillaume}
\tagline{Data Integration Specialist | Python, JavaScript/TypeScript, API, ETL, Data Quality \& Automation}

\photoL{2.8cm}{moi}

\personalinfo{%
  \email{guillaume.boileaupro@gmail.com}
  \phone{+33 6 59 94 85 84}
  \location{Alpes-Maritimes, FRANCE}
  \linkedin{guillaume-boileau-phd-891b81265}
  \researchgate{Guillaume-Boileau-2}
  \orcid{0000-0002-3576-6968}
}

\makecvheader
\vspace{-0.3cm}

\AtBeginEnvironment{itemize}{\small}
\columnratio{0.64}

\begin{paracol}{2}

\cvsection{Experience}

\cvevent{\textbf{Software Engineer - Data Flows, APIs \& Operational Integrations}}{VEV Platform Services France}{2025 -- 2026}{Valbonne, France}

\textbf{Context:} Software development and operational support for a CPMS platform connected to EV charging infrastructure, field equipment and external data exchanges.

\textbf{Objective:} Contribute to reliable operational data flows, system integrations and technical follow-up between software services, charging infrastructure and business operations.

\begin{itemize}
  \item \textbf{Data integration:} Implemented and monitored FTP-based operational data exchanges used for infrastructure follow-up and technical operations.
  \item \textbf{System interfaces:} Worked on software services connected to EV charging stations, operational equipment status and platform-side data flows.
  \item \textbf{Data consistency:} Checked consistency between CPMS services, external data exchanges, infrastructure behaviour and operational expectations.
  \item \textbf{Issue analysis:} Investigated production issues involving data exchanges, service behaviour, interface continuity and operational reliability.
  \item \textbf{Automation logic:} Supported automation-oriented workflows for operational data handling, monitoring and technical diagnostics.
  \item \textbf{Technical communication:} Structured technical observations and incident analyses to support prioritisation and decision-making.
  \item \textbf{Agile environment:} Used Zenhub for backlog follow-up, task tracking and coordination within an iterative software development environment.
  \item \textbf{Software stack:} Developed web and backend services using TypeScript, Angular and Node.js.
\end{itemize}

\divider

\cvevent{\textbf{Senior R\&D Engineer - Data Pipelines, Model Integration \& Validation}}{Laboratoire Lagrange, Observatoire de la Côte d'Azur, CNES}{2023 -- 2025}{Nice, France}

\textbf{Context:} Engineering and R\&D activities for the LISA ESA/NASA space mission, involving model integration, simulation workflows, data processing, validation campaigns and international coordination.

\textbf{Objective:} Design and maintain Python-based workflows to integrate heterogeneous models, structure technical data, validate outputs and support collaborative analysis.

\begin{itemize}
  \item \textbf{Data pipelines:} Developed Python workflows for model integration, data transformation, output comparison and validation.
  \item \textbf{Data management:} Structured simulation outputs, model configurations, technical repositories and analysis results to improve traceability and reuse.
  \item \textbf{Data quality:} Implemented consistency checks, comparison logic and validation rules to identify discrepancies between models and outputs.
  \item \textbf{Integration work:} Integrated heterogeneous external model results into a common analysis and review framework.
  \item \textbf{Technical documentation:} Used Confluence to organise documentation, coordination notes, validation follow-up and knowledge sharing.
  \item \textbf{Collaboration:} Worked with multidisciplinary contributors and translated complex technical results into clear synthesis material.
  \item \textbf{Process improvement:} Formalised reproducible workflows, comparison methods, validation checks and documentation practices.
  \item \textbf{Software platform:} Developed CATpyLISA, a Python platform for model integration, comparison, validation and technical review.
  \textcolor{DarkPastelRed}{\href{https://gitlab.oca.eu/gboileau/lisa_l3_tools}{\bf GitLab:CATpyLISA}}
\end{itemize}

\divider

\cvevent{\textbf{Data Scientist - Data Processing, Cleaning \& Technical Analysis}}{Universiteit Antwerpen, Belgium}{2021 -- 2023}{Antwerp, Belgium}

\textbf{Context:} Data analysis and performance studies for precision instruments in a demanding technical environment linked to gravitational-wave detector performance.

\textbf{Objective:} Build robust and reproducible data-analysis workflows to identify, characterise and mitigate sources affecting measurement quality.

\begin{itemize}
  \item \textbf{Python pipelines:} Developed Python-based workflows for data processing, statistical analysis and noise-source identification.
  \item \textbf{Data cleaning:} Processed, filtered and structured complex datasets to support reproducible technical analyses.
  \item \textbf{Data transformation:} Converted raw and derived data into exploitable formats for analysis, comparison and reporting.
  \item \textbf{Quality assessment:} Analysed data consistency, measurement stability and performance limitations in complex systems.
  \item \textbf{Reporting:} Produced traceable technical analyses and synthesis material for international technical stakeholders.
  \item \textbf{Problem solving:} Investigated environmental and instrumental constraints affecting system performance and data quality.
\end{itemize}

\divider

\cvevent{\textbf{Junior R\&D Engineer - Simulation, Data Analysis \& Automation}}{ARTEMIS, Observatoire de la Côte d'Azur}{2018 -- 2021}{Nice, France}

\textbf{Context:} R\&D activities for gravitational-wave mission preparation, involving signal analysis, simulation, technical documentation and noise characterisation.

\textbf{Objective:} Develop analysis methods, automate simulation workflows and document results in collaborative technical environments.

\begin{itemize}
  \item \textbf{Simulation workflows:} Built simulation approaches for complex signal environments, uncertainty studies and large-scale datasets.
  \item \textbf{Data analysis:} Developed methods for separating overlapping low-SNR signals in complex datasets.
  \item \textbf{Automation:} Implemented reproducible analysis workflows for repeated technical studies and comparison tasks.
  \item \textbf{Technical investigations:} Investigated noise sources through cross-sensor correlation analyses in diagnostic contexts.
  \item \textbf{Documentation:} Documented methods, assumptions, outputs and conclusions for collaborative review.
\end{itemize}

\switchcolumn

\cvsection{Profile for Mi4}

\textbf{Data integration-oriented engineer with strong experience in Python workflows, data pipelines, data quality, system interfaces, automation and technical communication}. Background developed through complex scientific and industrial environments, including \textbf{EV charging infrastructure software}, \textbf{LISA ESA/NASA}, \textbf{LIGO--Virgo--KAGRA} and \textbf{Einstein Telescope}.

For the \textbf{Data Integration Specialist} role, I bring a strong ability to analyse integration needs, structure data, automate transformations, monitor data flows, investigate anomalies and communicate clearly with both technical and non-technical stakeholders.

\begin{itemize}
  \item Strong experience in Python-based data processing, transformation, validation and automation workflows.
  \item Recent industrial software experience involving operational data flows, platform services and connected equipment.
  \item Practical experience with JavaScript/TypeScript, Node.js, Angular, Git, GitLab, Linux and Docker.
  \item Ability to work on APIs, external interfaces, data exchanges, system integration and operational reliability.
  \item Strong focus on data quality: consistency checks, traceability, reproducibility, validation evidence and anomaly investigation.
  \item Comfortable translating complex technical problems into clear explanations for multidisciplinary teams.
  \item Experience structuring technical repositories, documentation, workflows and data outputs for reuse and maintainability.
  \item Curious, autonomous and methodical profile, used to solving complex problems in evolving technical environments.
  \item Fast learner able to ramp up on CRM/ERP integration tools, HubSpot ecosystems, n8n and low-code automation platforms.
\end{itemize}

\cvsection{Languages}

\cvskill{English}{4}
\divider
\cvskill{Spanish}{2}
\divider

\medskip

\cvsection{Education}

\cvevent{PhD in Physics}{Université Côte d'Azur}{2018 -- 2021}{Nice, France}

\divider

\cvevent{Selective Graduate Program in Fundamental Physics (Magistère)}{Université Paris-Saclay}{2016 -- 2018}{Orsay, France}

\divider

\cvevent{MSc in Astronomy and Astrophysics}{Observatoire de Paris / Université Paris-Saclay}{2017 -- 2018}{Paris, France}

\divider

\cvevent{BSc in Fundamental Physics}{Université Paris-Saclay}{2015 -- 2016}{Orsay, France}

\cvsection{Professional Training}

\cvevent{Supervised Machine Learning (Regression \& Classification)}{DeepLearning.AI - Andrew Ng}{2020}{Coursera}

\end{paracol}

\cvsection{Projects}

\cvevent{CATpyLISA - Data Integration, Validation \& Automation Platform}{Python Platform for the LISA ESA/NASA Space Mission}{2023 -- 2025}{}

Development of a Python platform dedicated to large-scale model integration, data structuring, output comparison, validation and technical review for LISA mission-level workflows.

\textbf{Role:} Python development, data pipeline design, data structuring, validation logic, consistency checks, technical documentation and coordination with multiple contributors.

\textbf{Relevance for Mi4:} Experience directly transferable to data integration: heterogeneous data sources, automated transformations, traceability, validation, documentation and technical communication.

\divider

\cvevent{Operational Data Flows for EV Charging Infrastructure}{CPMS Platform - VEV Platform Services France}{2025 -- 2026}{}

Contribution to software services supporting operational data exchanges and connected infrastructure behaviour for EV charging systems.

\textbf{Role:} Development and follow-up of data exchange services, production issue analysis, operational diagnostics, interface reliability and technical coordination.

\textbf{Relevance for Mi4:} Practical industrial experience with operational data flows, external interfaces, technical troubleshooting and service-oriented software.

\divider

\cvevent{International Scientific and Technical Collaborations}{LISA, LIGO--Virgo--KAGRA and Einstein Telescope}{2018 -- now}{}

Contribution to collaborative developments in data analysis, signal characterisation, software workflows, validation activities and technical investigations within large international collaborations.

\textbf{Role:} Development of analysis tools, contribution to data-processing workflows, technical studies, documentation and collaborative review.

\textbf{Relevance for Mi4:} Ability to work with multidisciplinary teams, clarify complex technical issues, document methods and deliver reliable outputs in demanding collaborative environments.

\cvsection{Strengths}

\vspace{-.3cm}

\cvsubsection{Data Integration, ETL \& Automation}

{\small
\cvtag{Data Integration}
\cvtag{ETL Workflows}
\cvtag{Data Pipelines}
\cvtag{Data Cleaning}
\cvtag{Data Transformation}
\cvtag{Data Import}
\cvtag{Data Migration Logic}
\cvtag{Data Consolidation}
\cvtag{Workflow Automation}
\cvtag{Operational Data Flows}
\cvtag{System Integration}
\cvtag{Interface Monitoring}
\cvtag{Anomaly Investigation}
\cvtag{Process Improvement}
\cvtag{Low-Code Ramp-Up}
\cvtag{n8n Ramp-Up}
}

\divider\smallskip
\vspace{-.3cm}

\cvsubsection{APIs, Software \& Technical Tools}

{\small
\cvtag{Python}
\cvtag{JavaScript}
\cvtag{TypeScript}
\cvtag{Node.js}
\cvtag{Angular}
\cvtag{REST API}
\cvtag{Webhooks Logic}
\cvtag{FTP}
\cvtag{SQL}
\cvtag{Linux}
\cvtag{Bash}
\cvtag{Git}
\cvtag{GitLab}
\cvtag{GitHub}
\cvtag{Docker}
\cvtag{Confluence}
\cvtag{Zenhub}
\cvtag{AI-Assisted Coding}
\cvtag{Technical Debugging}
}

\divider\smallskip
\vspace{-.3cm}

\cvsubsection{Data Quality, CRM Logic \& Collaboration}

{\small
\cvtag{Data Quality}
\cvtag{Data Consistency}
\cvtag{Traceability}
\cvtag{Validation Checks}
\cvtag{Documentation}
\cvtag{Technical Synthesis}
\cvtag{Requirements Analysis}
\cvtag{Business Needs Translation}
\cvtag{CRM Integration Logic}
\cvtag{ERP Integration Logic}
\cvtag{HubSpot Ramp-Up}
\cvtag{Client-Oriented Mindset}
\cvtag{Cross-Functional Collaboration}
\cvtag{Problem Solving}
\cvtag{Autonomy}
\cvtag{Methodical Approach}
}

\end{document}